\chapter*{Portfolio Context and Evidence}
\thispagestyle{fancy}
\addcontentsline{toc}{chapter}{Portfolio Context and Evidence}

\noindent
\textbf{Project Context.}
This work originated from a graduate-level electromagnetic simulation exercise.
The present document is an independently prepared public portfolio edition intended
to demonstrate plane-wave boundary modeling, dielectric-slab impedance transformation,
frequency-domain simulation, parameter-sweep interpretation, and engineering
verification.

\vspace{0.5em}

\noindent
\textbf{Public Portfolio Edition.}
The original academic material has been reorganized into the same restrained
university-style architecture used by the other public engineering reports in this
portfolio. Student-identification information, course and instructor details,
institutional branding, and administrative submission language have been removed.
The technical design inputs and preserved simulation evidence are retained. This
document does not represent an official publication or endorsement by any university,
instructor, employer, or software vendor.

\vspace{0.5em}

\noindent
\textbf{Evidence Classification.}
Analytical evidence consists of the lossless transmission-line model for an
air--dielectric--air section, the half-wavelength matching condition, and an
independently executable Python script that regenerates the reflection magnitude and
wrapped phase. Simulated evidence consists of preserved CST Studio Suite Learning
Edition screenshots for geometry, boundary conditions, ports, parameter sweep,
fundamental-mode reflection coefficient, and electric- and magnetic-field results.
No fabricated hardware or laboratory measurement is involved.

\vspace{0.5em}

\noindent
\textbf{Simulation Scope.}
The public source archive contains exported figures and LaTeX reports but not the
native CST project, raw S-parameter exports, mesh statistics, or full solver logs.
The numerical conclusions are therefore limited to the preserved screenshots and
source-derived tables. The electric and magnetic side boundaries create a
TEM-like unit-cell environment; they do not constitute an infinite free-space
plane-wave model. A port-mode warning was reported during the sweep, and no formal
mesh- or port-convergence study was preserved.

\vfill

\begin{center}
    \textit{Prepared as part of the Hossein Molhem Engineering Portfolio.}
\end{center}
